\subsection{特征值的递减序列}

在本节中，我们假设 K = C。

我们置 \(\overline{\mathbf{N}} = \mathbf{N} \cup \{+\infty\} \subset \overline{\mathbf{R}}\) 。在本节中，若向量空间 E 不是有限维的，就说它有维数 \(+\infty \in \overline{\mathbf{N}}\) 。设 \(\mathrm{I_E} \subset \mathbf{N}\) 是 E 的满足 \(\mathrm{F} \neq \mathrm{E}\) 的有限维子空间 F 的维数所成的集。若 E 无穷维，我们有 \(\mathrm{I_E} = \mathbf{N}\) ，否则 \(\mathrm{I_E} = \{0, \ldots, \dim(\mathrm{E}) - 1\}\) 。当不致引起混淆时，我们置 \(\mathrm{I} = \mathrm{I_E}\) 。

设 \(u\) 是 E 的紧正（特别地，埃尔米特）自同态。\(u\) 的敏感谱是正实数的一个严格递减序列 \((\nu_{k})_{0\leqslant k<\mathrm{Card}(\mathrm{Sp}_{\mathrm{s}}(u))}\) 的值集（参见

III, p. 90, 命题 5）。对每个整数 \(k\) （满足 \(0 \leqslant k < \text{Card}(\text{Sp}_s(u))\) ），用 \(n_k\) 表示 \(\nu_k\) 的谱重数，其中 \(n_k \geqslant 1\) 。设 \(\text{M} \in \overline{\mathbf{N}}\) 是谱重数 \(n_k\) 之和；它是 \(u\) 的像的维数。我们有 \(\text{M} \leqslant \text{Card}(\text{I})\) 。

对 \(0 \leqslant n < M\) ，定义 \(\lambda_n(u) = \nu_k\) ，其中 \(k \geqslant 0\) 是唯一使下式成立的整数：

\[
n _ {0} + \dots + n _ {k - 1} \leqslant n <   n _ {0} + \dots + n _ {k}.
\]

置 \(\lambda_{n}(u)=0\) （若 \(n\in I\) 满足 \(n\geqslant M\) ）。这种情形只能在 I=N 且 \(\mathrm{Sp}_{\mathrm{s}}(u)\) 有限时发生（或者等价地，E 无穷维且 u 有限秩）。

序列 \((\lambda_{n}(u))_{n\in\mathrm{I}}\) 是递减的；对每个 \(\lambda\in\mathrm{Sp}_{\mathrm{s}}(u)\) ，使 \(\lambda_{n}(u)=\lambda\) 的整数 n 的个数等于 u 的特征值 \(\lambda\) 的谱重数。

我们约定俗成地说，\((\lambda_{n}(u))_{n\in\mathrm{J}}\) 是 u 的按重数重复的特征值递减序列。

命题 3. --- 设 u 是 E 的紧正自同态。存在 E 中的一个标准正交族 \((e_{n})_{n\in\mathrm{I}}\) ，使得对每个 \(x\in E\) ，有

\[
u (x) = \sum_ {n \in \mathrm{I}} \lambda_ {n} (u) \langle e _ {n} | x \rangle e _ {n}, \quad \langle x | u (x) \rangle = \sum_ {n \in \mathrm{I}} \lambda_ {n} (u) | \langle e _ {n} | x \rangle | ^ {2}.
\]

设 \(\mathrm{B} = (f_{j})_{j \in \mathrm{J}}\) 是 E 的使 u 为对角的标准正交基（IV, p. 150, 推论 1），\((\lambda_{j})_{j \in \mathrm{J}}\) 是 u 在基 B 中的特征值族。设 J' 是由 \(j \in J\) （使 \(\lambda_{j}\) 属于 u 的敏感谱）组成的集。

对每个 \(\lambda \in \mathrm{Sp}_{\mathrm{s}}(u)\) ，在整数 \(n\) （满足 \(0 \leqslant n < \mathrm{M}\) 与 \(\lambda_{n}(u) = \lambda\) ）与 \(j \in \mathrm{J}'\) （满足 \(\lambda_{j} = \lambda\) ）之间存在一个双射，因为这两个集有相同的基数，即 \(\lambda\) 的谱重数。对每个 \(\lambda\) 选取这样的双射，就定义了双射 \(\iota\) （从满足 \(0 \leqslant n < \mathrm{M}\) 的整数所成的集到集 \(\mathrm{J}'\) 的）。序列 \((e_n)_{0 \leqslant n < \mathrm{M}}\) 由置 \(e_n = f_{\iota(n)}\) （对 \(0 \leqslant n < \mathrm{M}\) ）在 E 中定义。它是 E 中的一个标准正交族。

在 M \textless{} Card(I) = +∞ 的情形，由 \(\{e_{0},\ldots,e_{M-1}\}\) 生成的空间 F 是有限维的，而它的正交 F° 是无穷维的；对 \((e_{n})_{n\geqslant M}\) 选取 F° 中的一个标准正交族。

对每个 \(x \in \mathrm{E}\) ，我们有

\[
u (x) = \sum_ {j \in J ^ {\prime}} \lambda_ {j} \langle f _ {j} \mid x \rangle f _ {j} = \sum_ {0 \leqslant n <   M} \lambda_ {n} (u) \langle e _ {n} \mid x \rangle e _ {n}
\]

因为 \(u(f_{j}) = 0\) （当 \(j \in J - J'\) ）。若 \(n \in I\) 满足 \(n \geqslant M\) ，我们有 \(\lambda_{n}(u) = 0\) ，由此得到命题的第一个公式。第二个公式由它得出。

命题 4. --- 设 u 是 E 的紧正自同态。设 \(f \in \mathcal{C}(\mathbf{R}_{+})\) 是连续递增函数，满足 \(f(0) = 0\) 。

自同态 \(f(u)\) 是紧且正的，且对每个 \(n \in \mathrm{I}_{\mathrm{E}}\) ，我们有 \(\lambda_n(f(u)) = f(\lambda_n(u))\) 。

由 III, p. 91, 命题 6, b) 与 I, p. 117, 命题 15, a)，自同态 \(f(u)\) 是紧且正的。\(f(u)\) 的谱是在 \(f\) 下 \(u\) 的谱的像（I, p. 111, 推论 2）。若 \(\lambda \in \mathrm{Sp}_{\mathrm{s}}(f(u))\) ，则 \(\lambda\) 的谱重数是 \(\mu \in \mathrm{Sp}(u)\) （满足 \(f(\mu) = \lambda\) ）的谱重数之和（III, p. 84, 推论 2）。由于 \(f\) 递增，序列 \((f(\lambda_n(u)))_{n \in \mathrm{I_E}}\) 是递减的。断言于是由序列 \((\lambda_n(f(u)))_{n \in \mathrm{I_E}}\) 的定义得出。

